% Generated draft: requires mathematical and source review.
\documentclass[11pt]{article}
\usepackage[margin=25mm]{geometry}
\usepackage{fontspec}
\setmainfont{Latin Modern Roman}
\usepackage{amsmath,amssymb,mathtools,unicode-math}
\setmathfont{Latin Modern Math}
\usepackage{enumitem,needspace,xurl,hyperref}
\setlength{\parindent}{0pt}
\setlength{\parskip}{5pt}
\setlength{\emergencystretch}{3em}
\newcommand{\N}{\mathbb{N}}
\newcommand{\Z}{\mathbb{Z}}
\newcommand{\Q}{\mathbb{Q}}
\newcommand{\R}{\mathbb{R}}
\newcommand{\C}{\mathbb{C}}
\newcommand{\F}{\mathbb{F}}
\newcommand{\A}{\mathbb{A}}
\newcommand{\PP}{\mathbb{P}}
\newcommand{\E}{\mathbb{E}}
\newcommand{\eps}{\varepsilon}
\newcommand{\dd}{\,\mathrm{d}}
\begin{document}
\section*{UnsolvedMath ID: 3100054}
\textbf{Source problem code:} catalogue entry
\subsection{Definitions and mathematical statement}
Interpret $p$ as a prime and let $\F_p=\Z/p\Z$. A transversal of the two coordinate parallel classes in $\F_p^2$ is a set $S$ meeting each horizontal and each vertical line in exactly one point. Equivalently,
\[
S=\{(x,f(x)):x\in\F_p\}
\]
for a permutation $f:\F_p\to\F_p$. An affine line is a set
\[
\ell_{a,b,c}=\{(x,y)\in\F_p^2:ax+by=c\},\qquad (a,b)\ne(0,0),
\]
and it is ordinary for $S$ if $|S\cap\ell_{a,b,c}|=2$. The literal source question, attributed to Solymosi, is whether every such transversal has an ordinary line. The description allows any two distinct parallel classes, which may be put into coordinate form by an invertible affine change of coordinates.

The literal statement lacks a non-collinearity hypothesis: for an odd prime, a transversal that is itself an affine line has no ordinary line. The supplied research summary refers to a ``corrected statement'' without stating it explicitly. Its natural nontrivial candidate is
\[
\forall p\text{ prime}\ \forall S\subseteq\F_p^2,\quad
\bigl(S\text{ is a transversal and is not contained in a line}\bigr)
\Longrightarrow\exists\ell:\ |S\cap\ell|=2.
\]
This candidate records the apparent intended correction rather than silently changing the supplied question. The case $p=2$ also needs care in the summary's counting claim: a full affine line then has exactly two points and is ordinary.
\subsection{Short English statement}
Must a non-collinear permutation graph over a prime field have a line through exactly two of its points? The source omits non-collinearity, while its research summary discusses a corrected conjecture.
\subsection{Sources}
\begin{itemize}
\item Local UnsolvedMath registry, permanent ID 3100054, exact source code catalogue entry, attributed to Solymosi; status partially solved, classification PARTIAL-PROGRESS.
\item Joshua Cooper, \emph{Combinatorial Problems I Like}, October 2020 snapshot, Combinatorial Number Theory, source-order bullet 54, \url{https://people.math.sc.edu/cooper/combprob.html}. Supplied HTML extraction accessed July 29, 2026. The snapshot labels the item open; later resolution was not checked item by item. The provenance flags status and redistribution permission NEEDS\_REVIEW; public access alone was verified.
\item The supplied research summary claims a new self-contained analytic argument for the corrected conjecture for $p\leq5$, and exhaustive verification through $p\leq13$, with the no-ordinary-line transversals described as the $p(p-1)$ affine lines. No argument, computation certificate, or publication reference is supplied, and the last description requires excluding $p=2$. The summary calls the problem open for $p\geq17$.
\item The summary names R\'edei's direction bound without a precise citation. For a hypothetical non-collinear no-ordinary-line transversal it reports at least $(p+3)/2$ determined slopes, at most $(p-1)/2$ determined slopes through each point, and the inequality $3|D|\leq\sum_{x\in S}r(x)\leq p(p-1)/2$. Here $D$ denotes determined directions and $r(x)$ is apparently their number through $x$; the source does not supply the promised preceding argument or detailed definitions. These are recorded claims, not verified inputs to the formulation.
\end{itemize}
\end{document}
